\section{Objectifs fonctionnels}

% ----------------------------------------------------------------
% 📌 À quoi sert cette section ?
% ----------------------------------------------------------------
% Cette partie liste de manière claire et synthétique les principales fonctionnalités attendues
% ou actions que le système/projet doit permettre à l’utilisateur d’accomplir.
% Il s’agit d’une vue fonctionnelle globale (le "quoi"), sans entrer dans les détails techniques (le "comment").

% ----------------------------------------------------------------
% 🧭 Comment la remplir ?
% ----------------------------------------------------------------
% ➤ Rédigez les objectifs fonctionnels sous forme de liste ou d’User Stories.
% ➤ Les User Stories suivent souvent la forme : 
%     "En tant que [rôle], je veux [fonctionnalité] afin de [bénéfice]."
% ➤ Classez-les par ordre logique ou de priorité si possible.
% ➤ Soyez clairs, concis, et évitez les termes trop techniques.

% ----------------------------------------------------------------
% 🧾 Exemple sous forme d’User Stories :
% ----------------------------------------------------------------
% \begin{itemize}
%   \item En tant qu’étudiant, je veux visualiser le résultat de mon code SVG en temps réel pour mieux comprendre mes erreurs.
%   \item En tant qu’utilisateur débutant, je veux accéder à une documentation simplifiée et interactive.
%   \item En tant que formateur, je veux pouvoir créer des défis personnalisés pour mes élèves.
% \end{itemize}

% ----------------------------------------------------------------
% 📊 Exemple avec tableau de priorités
% ----------------------------------------------------------------
% \begin{table}[H]
% \centering
% \begin{tabular}{|p{9cm}|c|}
% \hline
% Fonctionnalité (User Story) & Priorité \\
% \hline
% Visualiser du code SVG en temps réel & Haute \\
% Intégration de défis gamifiés & Moyenne \\
% Export des résultats & Basse \\
% \hline
% \end{tabular}
% \caption{Objectifs fonctionnels du projet}
% \end{table}

% ----------------------------------------------------------------
% ✍️ À vous de jouer : complétez ci-dessous
% ----------------------------------------------------------------

% (Départ suggéré)
Les objectifs fonctionnels du projet sont les suivants :

\begin{itemize}
  \item [À compléter]
\end{itemize}
